\renewcommand{\reportmonth}{Mei}
\renewcommand{\signaturedate}{2 Juni 2026}
\newpage
\TabPositions{4.5cm}

\begin{center}
    \textbf{LOG BOOK \\ MAGANG INDUSTRI}
\end{center}

\raggedright%
Program Studi           \tab: {\studyprogram} \\
Semester/Th. Akademik   \tab: {\semester} / {\academicperiod} \\
NIM/Nama Mahasiswa      \tab: {\authornim} / {\authorname} \\
Perusahaan              \tab: {\company} \\
Job Title               \tab: {\jobposition} \\
Departemen              \tab: {\jobdepartment} \\
Mentor                  \tab: {\mentor} \\
Bulan                   \tab: {\reportmonth} {\reportyear} \\
\vspace{12pt}

\begin{logtable}

    % =============================================================
    % Minggu Ke-17 (4 Mei -- 8 Mei 2026)
    % =============================================================

    \logentry{1.}
    {
        Senin \linebreak 4 Mei 2026 \linebreak
        08:00--17:35 \linebreak % chktex 8
        WFO
    }{
        \begin{loglist}
            \item Menyelesaikan Task \#13848 (IMPROVEMENT): membuat mekanisme navigasi
            \textit{split screen} antara React dan Toraja
            \item Melakukan \textit{merge} dan \textit{clean up} pekerjaan Task \#13410 yang
            sudah selesai di akhir April
        \end{loglist}
    }{
        \begin{logwork}
            \item (\#13848) Menambahkan fungsi \texttt{redirectSplitScreen} yang memungkinkan halaman React mengirim sinyal navigasi ke Toraja melalui \textit{iframe}
            \item (React) Menambahkan konfigurasi \textit{routing} Toraja ke aplikasi React, serta melakukan \textit{minor UI changes} dan \textit{merge branch master} ke branch 13410
            \item (Toraja) Menambahkan \texttt{bp\_id} ke \texttt{GetPilihNomorLotView}, menonaktifkan \textit{iframe} Riwayat Stok yang sudah dipindahkan, memperbaiki \textit{typo} \texttt{http} $\rightarrow$ \texttt{https}, dan melakukan beberapa \textit{merge branch master} ke branch 13410
        \end{logwork}
    }

    \logentry{2.}
    {
        Selasa \linebreak 5 Mei 2026 \linebreak
        07:58--17:21 \linebreak % chktex 8
        WFO
    }{
        \begin{loglist}
            \item Melakukan penyesuaian bahasa dan tampilan di beberapa komponen React setelah
            hasil \textit{code review} dari Mentor
        \end{loglist}
    }{
        \begin{logwork}
            \item (React) Mengubah \textit{date formatter} agar menggunakan bahasa Indonesia secara konsisten, termasuk pada komponen \textit{date range selector}
            \item (React) Menambahkan \textit{skeleton loader} baru (\texttt{listbased-skeleton}) yang tampil ketika proses pengambilan data dari \textit{backend} memakan waktu lebih lama
            \item (React) Menyembunyikan \textit{separator} di \texttt{dashboard-inventaris} saat tidak ada panel yang dibuka agar tampilan lebih bersih
        \end{logwork}
    }

    \logentry{3.}
    {
        Rabu \linebreak 6 Mei 2026 \linebreak
        08:02--17:48 \linebreak % chktex 8
        WFO
    }{
        \begin{loglist}
            \item Menyelesaikan Task \#14125 dan \#14130 (BUG): \textit{column mismatch} dan
            \textit{state handling} pada Stok Opname
            \item Mendapat Task \#14129 (BUG): \texttt{QuantityPicker} bertumpuk dengan
            \textit{scrollbar}
            \item Menyiapkan \textit{environment} baru untuk modul \textit{Stock Management} di
            Toraja dan packages
        \end{loglist}
    }{
        \begin{logwork}
            \item (\#14125 \& \#14130) Memperbaiki ketidaksesuaian nama kolom antara \textit{payload} \texttt{POST} dan \textit{response} API, serta meningkatkan logika \textit{error handling} dan manajemen \textit{state} pada halaman Stok Opname
            \item (\#14129) Menetapkan \texttt{zIndex: 0} pada \texttt{QuantityPicker} agar tidak bertumpuk dengan \textit{scrollbar} komponen di sekitarnya
            \item (React) Membangun komponen \texttt{error-toast} dan \texttt{show-toast}, mulai membangun halaman \texttt{dashboard-inventaris}, menyesuaikan warna utama ke \textit{shade} 4 sesuai warna asli HealthyOne, dan memformat tanggal panjang ke bahasa Indonesia
            \item (Toraja + packages) Menambahkan konfigurasi \textit{environment} baru khusus untuk modul \textit{Stock Management}
        \end{logwork}
    }

    \logentry{4.}
    {
        Kamis \linebreak 7 Mei 2026 \linebreak
        08:05--17:17 \linebreak % chktex 8
        WFO
    }{
        \begin{loglist}
            \item Melanjutkan perbaikan Stok Opname berdasarkan hasil \textit{testing} QA
        \end{loglist}
    }{
        \begin{logwork}
            \item (React) Memisahkan \textit{state} \texttt{search query} Stok Opname dari \texttt{bp\_id} agar keduanya tidak saling memengaruhi
            \item (React) Menyembunyikan sementara komponen \textit{filter} pada \texttt{ListBased} Stok Opname yang belum berfungsi penuh agar tidak membingungkan pengguna
        \end{logwork}
    }

    \logentry{5.}
    {
        Jumat \linebreak 8 Mei 2026 \linebreak
        07:55--17:10 \linebreak % chktex 8
        WFO
    }{
        \begin{loglist}
            \item Menyelesaikan Task \#13848: memperbaiki nilai \texttt{is\_split\_screen} yang
            salah saat melakukan \textit{redirect}
        \end{loglist}
    }{
        \begin{logwork}
            \item (\#13848) Memperbaiki bug di mana nilai \texttt{is\_split\_screen} dikirim sebagai \texttt{false} saat seharusnya \texttt{true} ketika halaman dibuka dalam mode \textit{split screen} dari Toraja
        \end{logwork}
    }

    % =============================================================
    % Minggu Ke-18 (11 Mei -- 15 Mei 2026)
    % =============================================================

    \logentry{6.}
    {
        Senin \linebreak 11 Mei 2026 \linebreak
        07:59--17:33 \linebreak % chktex 8
        WFO
    }{
        \begin{loglist}
            \item Menyelesaikan Task \#14201 dan \#14281 (BUG): masalah pemilihan baris pada
            \texttt{PilihNomorLot}
            \item Menyelesaikan Task \#14282 (USER STORY): tampilkan \textit{error toast} jika
            item dipilih tanpa perubahan kuantitas
        \end{loglist}
    }{
        \begin{logwork}
            \item (React) Memperbaiki logika \textit{select row} berdasarkan ID dan ketidaksesuaian tipe data di halaman \texttt{pilih-nomor-lot}
            \item (\#14281 \& \#14201) Mengubah \textit{state} \texttt{selectedIds} dan menambahkan \texttt{allLoadedRows} untuk menyimpan baris dari seluruh halaman agar seleksi lintas halaman berjalan dengan benar
            \item (\#14282) Menampilkan \textit{error toast} ketika pengguna memilih item tetapi tidak mengubah kuantitas sebelum menyimpan
            \item (React) Meningkatkan UI \texttt{ListBased}: memperbaiki \textit{state} \texttt{viewType}, menambahkan tombol \textit{extra} untuk tampilan \textit{card}, serta memperbaiki tinggi baris dan tampilan \textit{pagination}
        \end{logwork}
    }

    \logentry{7.}
    {
        Selasa \linebreak 12 Mei 2026 \linebreak
        08:03--17:40 \linebreak % chktex 8
        WFO
    }{
        \begin{loglist}
            \item Membangun komponen generik baru \textit{DataView} sebagai pengganti
            \textit{ListBased} yang lebih \textit{powerful} dan fleksibel
        \end{loglist}
    }{
        \begin{logwork}
            \item (React) Merancang dan membangun struktur komponen \textit{DataView} (\textit{ListBased v2}) beserta arsitektur internalnya: \texttt{TableView}, \texttt{CardView}, \textit{filter menu}, \textit{group by}, dan \textit{visible columns}
            \item (React) Mengimplementasikan sistem \textit{filtering} berbasis \textit{backend} pada \textit{DataView}, termasuk pengiriman parameter \textit{filter} ke API secara dinamis
        \end{logwork}
    }

    \logentry{8.}
    {
        Rabu \linebreak 13 Mei 2026 \linebreak
        08:07--17:50 \linebreak % chktex 8
        WFO
    }{
        \begin{loglist}
            \item Melanjutkan pengembangan komponen \textit{DataView} dengan menambahkan fitur
            \textit{card}, \textit{group by}, dan \textit{visible columns}
            \item Melakukan \textit{refactor} halaman \texttt{laporan-stok-opname} agar
            menggunakan \textit{DataView}
        \end{loglist}
    }{
        \begin{logwork}
            \item (React) Menambahkan fungsi \texttt{setCardHeader} pada \textit{DataView} untuk menetapkan kolom tertentu sebagai \textit{header card}
            \item (React) Melakukan \textit{refactor} halaman \texttt{laporan-stok-opname}: mengganti \texttt{ListBased} dengan \textit{DataView} beserta \textit{pagination}, \textit{filtering}, \textit{group by}, dan \textit{date select} yang sepenuhnya fungsional
            \item (React) Membuat fungsi \texttt{generateColumnDef} yang secara otomatis membuat definisi kolom dari \texttt{response.data} API
            \item (React) Menambahkan tampilan \textit{filter/group by} yang sedang aktif di bawah \textit{control bar} serta mengimplementasikan \textit{backend group by} dan fitur \textit{visible columns} (sembunyikan/tampilkan kolom)
        \end{logwork}
    }

    \logentry{9.}
    {
        Jumat \linebreak 15 Mei 2026 \linebreak
        08:00--17:28 \linebreak % chktex 8
        WFO
    }{
        \begin{loglist}
            \item Melengkapi fitur \textit{filtering} pada \textit{DataView} dan melakukan
            \textit{refactor} halaman \texttt{laporan-pembelian}
        \end{loglist}
    }{
        \begin{logwork}
            \item (React) Menambahkan fungsi \texttt{hideFilters} pada \textit{DataView} untuk menyembunyikan \textit{filter} tertentu dari \texttt{FilterMenu}, serta menambahkan tipe \textit{filter}: \textit{single select} dan \textit{multi select}
            \item (React) Melakukan \textit{refactor} halaman \texttt{laporan-pembelian}: mengganti \texttt{ListBased} dengan \textit{DataView} serta mengimplementasikan \textit{pagination}, \textit{filtering}, \textit{group by}, dan \textit{date range select} yang sepenuhnya fungsional
        \end{logwork}
    }

    \logevaluate{Dokumentasi Hasil Pekerjaan: Hasil Pekerjaan: Komponen DataView dengan fitur filter, group by, dan visible columns}{
        \noindent\includegraphics[width=\linewidth]{3.4-dataview-component.png}
    }

    % =============================================================
    % Minggu Ke-19 (18 Mei -- 22 Mei 2026)
    % =============================================================

    \logentry{10.}
    {
        Senin \linebreak 18 Mei 2026 \linebreak
        08:04--17:42 \linebreak % chktex 8
        WFO
    }{
        \begin{loglist}
            \item Mengembangkan fitur \textit{filtering} dan \textit{labeling} pada
            \texttt{dashboard-inventaris} dan \textit{DataView}
        \end{loglist}
    }{
        \begin{logwork}
            \item (React) Memperbaiki \textit{DataView}: menangani \textit{input} kosong pada \textit{multi search} dan memindahkan definisi kolom (\textit{coldef}) \texttt{dashboard-inventaris} ke file tersendiri
            \item (React) Membuat \textit{interface}, logika, dan komponen \textit{filtering} untuk \texttt{dashboard-inventaris} serta menambahkan kolom-kolom baru
            \item (React) Menambahkan dukungan untuk label dan \textit{placeholder} terpetakan dari halaman ke \textit{DataView} agar nama kolom dan \textit{filter} yang tampil bisa disesuaikan
            \item (React) Melakukan \textit{refactor} \texttt{laporan-pembelian} menggunakan \textit{DataView} sepenuhnya serta menambahkan tampilan \textit{mobile responsive} untuk \texttt{dashboard-inventaris}
        \end{logwork}
    }

    \logentry{11.}
    {
        Selasa \linebreak 19 Mei 2026 \linebreak
        07:57--17:35 \linebreak % chktex 8
        WFO
    }{
        \begin{loglist}
            \item Menambahkan fungsionalitas pembayaran di \texttt{laporan-pembelian} dan fitur
            \textit{selectable row} di \textit{DataView}
        \end{loglist}
    }{
        \begin{logwork}
            \item (React) Menambahkan fungsi \texttt{PutPurchaseReportManager} (tombol Bayar) di \texttt{laporan-pembelian} untuk memproses pembayaran laporan pembelian
            \item (React) Menambahkan \textit{overlay} pesan \textit{unauthorized error} pada \texttt{dashboard-inventaris}
            \item (React) Menambahkan fitur \textit{selectable row} pada \textit{DataView}, \textit{popover hover} pada \textit{header} kolom, dan \textit{skeleton loading} yang lebih baik
            \item (React) Memperbarui fungsi \texttt{formatIDR} dan menambahkan \texttt{parseError} ke utils
        \end{logwork}
    }

    \logentry{12.}
    {
        Rabu \linebreak 20 Mei 2026 \linebreak
        08:01--17:29 \linebreak % chktex 8
        WFO
    }{
        \begin{loglist}
            \item Menyelesaikan Task \#14486 (BUG): kuantitas tidak sesuai unit terkecil di
            \texttt{pilih-nomor-lot}
            \item Melakukan \textit{cleanup} dan berbagai perbaikan \textit{minor}
        \end{loglist}
    }{
        \begin{logwork}
            \item (\#14486) Memperbaiki bug \texttt{setQuantities} di \texttt{pilih-nomor-lot}: menambahkan logika konversi ke unit terkecil di dalam \texttt{useEffect} agar \texttt{Pick Unit Select} tidak mengalami ketidaksesuaian saat pertama kali dimuat
            \item (React) Melakukan \textit{cleanup} kode dan memperbarui package \texttt{Mantine} ke versi terbaru
            \item (React) Menambahkan \textit{label map} untuk nama kolom saat melakukan \textit{group by} di \textit{DataView} dan \texttt{dashboard-inventaris}
            \item (React) Memperbaiki \texttt{laporan-pembelian}: menambahkan \texttt{payment\_status\_label} untuk menentukan status tombol Bayar, menonaktifkan beberapa kolom yang belum relevan, dan memperbaiki \textit{field-column name mismatch}
            \item (React) Memperbaiki \textit{DataView}: mengganti \textit{search chip} menjadi \texttt{renderPill} agar proses hapus \textit{search} lebih responsif, serta memperbaiki nama kolom yang dipilih untuk \textit{group by}
        \end{logwork}
    }

    \logentry{13.}
    {
        Kamis \linebreak 21 Mei 2026 \linebreak
        08:08--17:45 \linebreak % chktex 8
        WFO
    }{
        \begin{loglist}
            \item Menyelesaikan Task \#14263 (BUG): ketidakkonsistenan tinggi widget dan ukuran
            ikon pada Flutter
            \item Menyelesaikan Task \#13999 (BUILD): menghubungkan \texttt{dashboard-inventaris}
            React dengan Toraja \textit{backend}
            \item Menyelesaikan Task \#14553 (BUG) di repository
            \texttt{healthcare-productdetection}
        \end{loglist}
    }{
        \begin{logwork}
            \item (\#14263) Memperbaiki ketidakkonsistenan tinggi widget dan ukuran ikon pada \texttt{TextForm}, \texttt{TimePicker}, \texttt{SingleDropdown}, dan \texttt{MultiDropdown} yang terkait sistem \textit{font size} tabel
            \item (React) Menghapus file konfigurasi \texttt{pnpm} yang tidak digunakan serta memperbaiki warna \textit{hover} dan \textit{selected} pada \textit{list-based}
            \item (\#13999) Membangun integrasi \textit{backend} untuk \texttt{dashboard-inventaris}: menghubungkan \texttt{laporan-pembelian} dan \texttt{laporan-stok-opname} React ke Toraja menggunakan \textit{iframe} dan API
            \item (\#14553) Menghapus \textit{flag} \texttt{is\_currency} pada kolom 'Alasan Pengajuan PS/SP' yang menyebabkan format data tidak benar
        \end{logwork}
    }

    \logentry{14.}
    {
        Jumat \linebreak 22 Mei 2026 \linebreak
        07:53--17:21 \linebreak % chktex 8
        WFO
    }{
        \begin{loglist}
            \item Menyelesaikan Task \#14579 (BUG): perubahan versi \texttt{Mantine} menyebabkan
            \textit{checkbox} di \textit{DataView} tidak berfungsi
            \item Menambahkan \textit{font size provider} pada \textit{DataView}
        \end{loglist}
    }{
        \begin{logwork}
            \item (React, \#14579) Memperbaiki \textit{module version change} yang menyebabkan \textit{checkbox} \texttt{Mantine} tidak ter-\textit{render} dengan benar di \textit{DataView}
            \item (React) Menambahkan kemampuan komponen \textit{custom render} di dalam \textit{DataView} untuk mengikuti \textit{font size provider} yang aktif
            \item (React) Menonaktifkan pengurutan untuk kolom yang memiliki komponen \textit{render} kustom di \texttt{pilih-nomor-lot}
            \item (React) Mengoptimalkan penanganan \texttt{laporanType} di \texttt{laporan-pembelian} agar data tidak hilang setelah proses \textit{fetch}
        \end{logwork}
    }

    % =============================================================
    % Minggu Ke-20 (25 Mei -- 29 Mei 2026)
    % =============================================================

    \logentry{15.}
    {
        Senin \linebreak 25 Mei 2026 \linebreak
        07:58--17:36 \linebreak % chktex 8
        WFO
    }{
        \begin{loglist}
            \item Melengkapi fungsionalitas \textit{DataView} dengan fitur \textit{selectable}
            yang sepenuhnya berjalan di \texttt{TableView} dan \texttt{CardView}, serta
            \textit{custom controls}
        \end{loglist}
    }{
        \begin{logwork}
            \item (React) Menambahkan fitur \textit{selectable} penuh pada \texttt{TableView} dan \texttt{CardView}: komponen \textit{render} dinamis mengikuti \texttt{TableFontSize}, memperbaiki algoritma pengurutan, dan melakukan various perbaikan \textit{styling}
            \item (React) Menambahkan \textit{custom controls} untuk menyembunyikan komponen tertentu dari \textit{DataView} (seperti \texttt{FilterMenu}, \texttt{SearchBar}, \texttt{Select ViewType}), menambahkan tampilan \textit{empty state}, serta mengoptimalkan performa komponen
            \item (React) Menambahkan \textit{transition} pada komponen \texttt{footer-button}
            \item (React) Menambahkan \texttt{initialMount} \textit{ref} untuk mengatasi bug di mana kolom \textit{visible} menghilang ketika \textit{DataView} melakukan \textit{rerender}
        \end{logwork}
    }

    \logentry{16.}
    {
        Selasa \linebreak 26 Mei 2026 \linebreak
        08:00--17:15 \linebreak % chktex 8
        WFO
    }{
        \begin{loglist}
            \item Melakukan \textit{internal review} hasil pekerjaan \textit{DataView} bersama
            Mentor dan tim QA
        \end{loglist}
    }{
        \begin{logwork}
            \item Mempresentasikan hasil fitur \textit{DataView} terbaru kepada Mentor untuk
            mendapatkan \textit{feedback} dan arah pengembangan selanjutnya
            \item Mendiskusikan prioritas perbaikan dan pengembangan bersama tim QA berdasarkan
            hasil \textit{testing} komponen \textit{DataView}
        \end{logwork}
    }

    \logevaluate{Dokumentasi Hasil Pekerjaan: Hasil Pekerjaan: Layar Detail Klinik Canvasser}{
        \noindent\includegraphics[width=\linewidth]{3.4-canvasser-detailklinik.png}
    }

    \logentry{17.}
    {
        Rabu \linebreak 27 Mei 2026 \linebreak
        08:03--17:22 \linebreak % chktex 8
        WFO
    }{
        \begin{loglist}
            \item Merencanakan dan mulai mempersiapkan fitur \textit{export} data dan
            \textit{multiple group by} pada \textit{DataView} dan
            \texttt{dashboard-inventaris}
        \end{loglist}
    }{
        \begin{logwork}
            \item Melakukan riset implementasi \textit{export} ke CSV dan XLSX menggunakan
            \textit{library} yang kompatibel dengan \textit{stack} yang digunakan
            \item Merancang arsitektur untuk fitur \textit{multiple group by keys} pada
            \textit{DataView}
        \end{logwork}
    }

    \logentry{18.}
    {
        Kamis \linebreak 28 Mei 2026 \linebreak
        07:51--17:44 \linebreak % chktex 8
        WFO
    }{
        \begin{loglist}
            \item Melakukan \textit{refactor} besar pada \texttt{dashboard-inventaris} dan
            \texttt{pilih-nomor-lot} serta berbagai perbaikan performa pada
            \textit{DataView}
        \end{loglist}
    }{
        \begin{logwork}
            \item (React) Mengoptimalkan \texttt{dashboard-inventaris}: memecah komponen \texttt{laporan-pembelian} menjadi 3 komponen dan \textit{service} terpisah untuk menghindari \textit{data racing}, mengoptimalkan \textit{styling} dan logika \texttt{laporan-stok-opname}, serta menambahkan \texttt{TableFontSize} yang dinamis
            \item (React) Mengoptimalkan \texttt{pilih-nomor-lot}: mengganti \texttt{ListBased} dengan \textit{DataView}, mengoptimalkan parameter \textit{service}, serta mempertahankan \textit{state} \texttt{selectedIds} saat \textit{pagination} atau \textit{search} dilakukan
            \item (React) Memperbaiki \textit{DataView}: menambahkan \textit{overflow wrap} saat \textit{cardview} terlalu sempit, memperbaiki \textit{filtering} sebelum data di-\textit{render} oleh \texttt{AgGrid}, dan menonaktifkan sementara \textit{cell align}
            \item (React) Memperbaiki \textit{state} \textit{group by} menjadi \texttt{null} jika kosong, menambahkan tombol hapus pada \textit{chip filter} tipe \textit{group}, dan memperbaiki \textit{double search} yang dipicu oleh \texttt{onBlur}
            \item (React) Menambahkan \textit{advanced styling functionality} pada \texttt{quantity-picker} serta memperbaiki \texttt{formatIDR} agar menangani nilai \texttt{NaN}
        \end{logwork}
    }

    \logentry{19.}
    {
        Jumat \linebreak 29 Mei 2026 \linebreak
        08:05--17:32 \linebreak % chktex 8
        WFO
    }{
        \begin{loglist}
            \item Menyelesaikan beberapa \textit{bug fix} Flutter: Task \#13670, \#14699, dan
            perbaikan \texttt{text-form}
            \item Menyelesaikan Task \#14696 (IMPROVEMENT): fitur \textit{export} CSV dari
            \textit{DataView} dan \texttt{dashboard-inventaris}
            \item Menyelesaikan Task \#14697 (IMPROVEMENT): fitur \textit{multiple group by keys}
            pada \textit{DataView} dan \texttt{dashboard-inventaris}
        \end{loglist}
    }{
        \begin{logwork}
            \item (\#13670) Menonaktifkan fitur \textit{group by} untuk tipe tampilan \textit{card-vertical} pada \texttt{ListBased} yang belum mendukung fitur tersebut
            \item (\#14699) Membungkus komponen \textit{error message} dengan \texttt{Flexible} agar tidak \textit{overflow} terhadap \textit{widget} lain di sekitarnya
            \item Memperbaiki \texttt{text-form}: menambahkan \textit{ternary} untuk menentukan
            ukuran \textit{font} berdasarkan \texttt{isForEditTable} serta memperbaiki
            warna simbol \textit{required}
            \item (\#14696) Menambahkan fungsi \textit{export} CSV di sisi \textit{frontend} pada \textit{DataView} dan halaman \texttt{dashboard-inventaris}
            \item (\#14697) Menambahkan fitur \textit{multiple group by keys} pada \textit{DataView} dan \texttt{dashboard-inventaris}, memungkinkan data dikelompokkan berdasarkan lebih dari satu kolom sekaligus
        \end{logwork}
    }

    % =============================================================
    % Evaluasi Bulanan
    % ============================================================= 
    \logevaluate{Analisis dan Pembahasan}{
        Di awal bulan, pekerjaan berfokus pada penyelesaian sisa \textit{alignment} dari April (Task \#13286, \#13592, \#13999, \#14000, \#13848), yang kemudian langsung diikuti oleh serangkaian \textit{bug fix} dan \textit{user story} pada halaman Stok Opname (\#14125--\#14133, \#14201, \#14176).

        Development komponen generik \texttt{DataView} (Task \#14855) pada minggu kedua
        merupakan tugas terbesar selama menjalani magang karena fungsionalitasnya yang
        kompleks. Komponen ini dirancang sebagai pengganti \texttt{ListBased} yang
        lebih canggih: mendukung \textit{pagination}, \textit{filtering}, \textit{group
            by}, \textit{visible columns}, \textit{selectable rows}, dan \textit{custom
            render components} yang semuanya bisa dikonfigurasi dari halaman pemakai. Dalam
        waktu dua minggu, \texttt{DataView} sudah diterapkan di
        \texttt{laporan-stok-opname}, \texttt{laporan-pembelian}, dan
        \texttt{dashboard-inventaris}.

        Di minggu ketiga dan keempat, pekerjaan terbagi antara fitur-fitur baru (Task
        \#14481 --- Layar Detail Klinik Canvasser, Task \#14586 --- migrasi komponen ke
        DataView, Task \#14696/\#14697 --- export dan multiple group by) serta berbagai
        \textit{bug fix} lintas repository (\#14271, \#14486, \#14553, \#14572,
        \#14579, \#14646, \#14685, \#14699). Pada periode ini peran \textit{full-stack}
        semakin terasa dengan tetap terlibat di sisi Toraja untuk keperluan endpoint
        Dashboard Inventaris. }

    \logevaluate{Capaian yang Telah Diselesaikan}{
        \begin{itemize}[noitemsep, topsep=2pt, leftmargin=12pt, label=\textbullet]
            \item Menyelesaikan Task \#13286 \& \#13592: halaman laporan Stok Opname dan laporan
                  Pembelian React (STAGING).
            \item Menyelesaikan Task \#13999: \textit{alignment} Dashboard Inventaris ke
                  \textit{backend} Toraja.
            \item Menyelesaikan Task \#14000: Riwayat Stok align ke \textit{backend}.
            \item Menyelesaikan Task \#13848: navigasi \textit{split screen} React--Toraja
                  berfungsi penuh.
            \item Menyelesaikan Task \#14125, \#14128--\#14131, \#14133, \#14176, \#14201:
                  serangkaian perbaikan halaman Stok Opname.
            \item Menyelesaikan Task \#14234: filter Dashboard Inventaris.
            \item Menyelesaikan Task \#14263: konsistensi ukuran widget pada \textit{editable
                      table} Flutter.
            \item Menyelesaikan Task \#14271: bug \textit{double entry} supplier di detail PSSP.
            \item Menyelesaikan Task \#14281 \& \#14282: perbaikan Stok Opname multi-halaman dan
                  validasi seleksi produk.
            \item Menyelesaikan Task \#14486: konversi kuantitas ke unit terkecil di
                  \texttt{pilih-nomor-lot}.
            \item Menyelesaikan Task \#14481: layar Detail Klinik Canvasser.
            \item Menyelesaikan Task \#14553: perbaikan tampilan alasan under branded Canvasser.
            \item Menyelesaikan Task \#14572 \& \#14579: perbaikan data tabel dan checkbox di
                  \texttt{laporan-pembelian}.
            \item Menyelesaikan Task \#14586: migrasi komponen \texttt{ListBased} ke
                  \texttt{DataView} di seluruh halaman terdampak.
            \item Menyelesaikan Task \#14627: tombol Bayar tetap di halaman yang sama.
            \item Menyelesaikan Task \#14127: perbaikan lanjutan \textit{font size} widget
                  Flutter.
            \item Menyelesaikan Task \#14646 \& \#14685: bug fix pasca-migrasi pada Dashboard
                  Inventaris dan Stok Opname.
            \item Menyelesaikan Task \#14699: perbaikan pesan validasi \textit{birthdate}
                  Flutter.
            \item Menyelesaikan Task \#13670: perbaikan \textit{Group By} filter pada
                  \texttt{ListBased} Flutter.
            \item Memulai Task \#14696, \#14697, \#14740, dan \#14374 yang akan dilanjutkan di
                  Juni.
        \end{itemize}
    }

    \logevaluate{Capaian yang Belum Diselesaikan}{
        \begin{itemize}[noitemsep, topsep=2pt, leftmargin=12pt, label=\textbullet]
            \item Task \#14696 (\textit{export} XLSX), \#14697 (\textit{multiple group by}),
                  \#14740 (\textit{fix sort jumlah data}), dan \#14374 (\textit{Diagnosis
                      Sekunder hilang}) sudah dimulai di akhir bulan dan akan diselesaikan di awal
                  Juni.
        \end{itemize}
    }

    \logevaluate{Kendala}{
        Pembangunan komponen \texttt{DataView} dari nol sambil sekaligus menggunakannya di beberapa halaman berbeda menimbulkan banyak bug yang tidak terduga --- terutama terkait \textit{state management} saat \textit{rerendering}, konflik \textit{scrollbar}, dan ketidaksesuaian tipe data. Debugging membutuhkan waktu cukup lama karena penyebabnya seringkali berasal dari interaksi antar komponen yang kompleks.
    }

    \logevaluate{Solusi}{
        Pembangunan \texttt{DataView} dilakukan secara bertahap dengan mengutamakan pengujian di setiap penambahan fitur baru. Ketika bug muncul, \textit{state flow} ditelusuri satu per satu menggunakan \texttt{console.log} sementara dan \textit{React DevTools}. Migrasi dari \texttt{ListBased} ke \texttt{DataView} dilakukan halaman per halaman (bukan sekaligus) untuk memastikan setiap halaman berfungsi sebelum lanjut ke berikutnya.
    }

\end{logtable}
\begin{flushright}
    {\city}, {\signaturedate}\\
    Mengetahui,\\
    Mentor {\jenisMBKM}\\[2cm]
    {\mentor}\\
\end{flushright}